\subsection{Resultats}
\label{sec:DSCresultats}

Lors de la chauffe rapide Fig. \ref{fig:chauffe_rapide}, nous constatons une temperature de transition vitreuse autour de 82 $^\circ$C et pic exo/endo-thermique de -35.58 J/g. (fusion==endo)

En revanche, le refroidissement rapide ne laisse pas au materiau le temps de correctement se restructurer. (qu'est ce que cela implique?)

\begin{flushleft}
\begin{minipage}{0.49\textwidth} 
\includegraphics[width=\textwidth]{dsc_chauffe.pdf}%
	\captionof{figure}{$\tan\delta$(T[$^\circ$C]) chauffe rapide}
       \label{fig:chauffe_rapide}
\end{minipage}
\hfill
\begin{minipage}{0.49\textwidth}
    \includegraphics[width=\linewidth]{dsc_amorphe.pdf}
	\captionof{figure}{$\tan\delta$(T[$^\circ$C]) refroidissement rapide}
        \label{fig:refroidissement_rapide}
\end{minipage}
\end{flushleft}
\\\\
Neanmoins, le comportement de l'echantillon change fortement quand nous lui permettons de se refroidir ou chauffer plus lentement (visible sur Figs. \ref{fig:chauffe_lente} et \ref{fig:refroidissement_lent}).

\begin{flushleft}
\begin{minipage}{0.49\textwidth}
    \includegraphics[width=\linewidth]{dsc_cristallisation_chaude.pdf}
        \captionof{figure}{$\tan\delta$(T[$^\circ$C]) chauffe lente}
        \label{fig:chauffe_lente}
\end{minipage}
\hfill
\begin{minipage}{0.49\textwidth}
    \includegraphics[width=\linewidth]{dsc_cristallisation_froide.pdf}
        \captionof{figure}{$\tan\delta$(T[$^\circ$C]) refroidissement lent}
        \label{fig:refroidissement_lent}
\end{minipage}
\end{flushleft}

\\\\
Quand nous chauffons l'echantillon plus lentement (voir Fig. \ref{fig:chauffe_lente}), nous observons sa temperature de transition vitreuse autour de 77 $^\circ$C ainsi qu'un pic endothermique vers 133.1 $^\circ$C et un pic exothermique autour de 250 $^\circ$C.
DEVELOPPER...

\\\\
Alors que le refroidir plus lentement (voir Fig. \ref{fig:refroidissement_lent}), nous laisse entrevoir plus clairement un pic. DEVELOPPER...
